\documentclass[10pt]{article}
\usepackage[T1]{fontenc}
\usepackage[utf8]{inputenc}
\usepackage[english]{babel}
\usepackage{lmodern}
\usepackage{amsmath,amssymb,mathtools}
\usepackage{enumitem}
\usepackage{tikz-cd}
\usepackage{microtype}
\usepackage[paperwidth=461pt,paperheight=684pt,top=34pt,bottom=34pt,left=38pt,right=38pt]{geometry}
\pagestyle{empty}
\setlength{\parindent}{0pt}
\setlength{\parskip}{2.5pt}
\setlist{nosep,leftmargin=2em}
\setlength{\emergencystretch}{2em}
\renewcommand{\thefootnote}{\fnsymbol{footnote}}
\newcommand{\UX}{\underline{X}}
\newcommand{\UL}{\underline{L}}
\newcommand{\US}{\underline{S}}
\newcommand{\cris}{\mathrm{cris}}
\newcommand{\DR}{\mathrm{DR}}
\pdfinfoomitdate=1
\pdfsuppressptexinfo=15
\pdftrailerid{<746961B93EFDEAEB4C4F7B8E6D5020C2><746961B93EFDEAEB4C4F7B8E6D5020C2>}
\pdfinfo{/Title (Lifting K3 surfaces to characteristic zero - English translation) /Author (P. Deligne; written up by Luc Illusie) /Subject (Translation from frozen diplomatic French source pages) /Producer (D039 deterministic normalized-edition builder) /CreationDate (D:19800101000000Z) /ModDate (D:19800101000000Z)}
\begin{document}
\setcounter{footnote}{0}
\fontsize{10.2pt}{12.04pt}\selectfont

\begin{center}
\textbf{EXPOSÉ IV}\\[2.5em]
\textbf{LIFTING K3 SURFACES TO CHARACTERISTIC ZERO}\\[2.0em]
by P. DELIGNE\\[0.4em]
(written up by Luc ILLUSIE)\footnote{Research Team associated with the C.N.R.S., no.~653.}
\end{center}

Rudakov and Shafarevitch [14] (see also Nygaard [12]) established that K3 surfaces admit no nonzero vector fields.

The purpose of this exposé is to prove, from this, that every polarized K3 surface in characteristic $p>0$ lifts to characteristic zero (1.8). The proof rests on the fact that the Rudakov--Shafarevitch result entails, for K3 surfaces, the existence of a good deformation theory, making it possible in particular to ``control'' the locus in the versal formal space along which a given nontrivial invertible sheaf deforms. The only ingredients used are, on the one hand, the usual relations in de Rham cohomology between the Gauss--Manin connection and the Kodaira--Spencer operation, and, on the other hand, the existence of a crystalline Chern class for invertible sheaves.

We denote by $k$ an algebraically closed field of characteristic $p>0$, and write $W=W(k)$ for the ring of Witt vectors over $k$.
\newpage
\setcounter{footnote}{0}
\fontsize{10.2pt}{12.04pt}\selectfont

\section*{1. STATEMENT OF THE LIFTING THEOREM}
In this section, $X_0$ denotes a K3 surface over $k$.

\noindent\textbf{PROPOSITION 1.1.}
\begin{enumerate}[label=\alph*)]
\item The Hodge spectral sequence
\[
E_1^{ij}=H^j(X_0,\Omega^i_{X_0/k})\Longrightarrow H^*_{\mathrm{DR}}(X_0/k)
\]
degenerates at $E_1$, and the matrix of Hodge numbers
$h^{ij}=\dim_k H^j(X_0,\Omega^i_{X_0/k})$ is
\[
\begin{matrix}1&0&1\\0&20&0\\1&0&1\end{matrix}.
\]
\item Let $T=T_{X_0/k}=(\Omega^1_{X_0/k})^\vee$ be the tangent bundle of $X_0$. Then $H^i(X_0,T)=0$ for $i=0$ or $2$, and $\dim_kH^1(X_0,T)=20$.
\item The crystalline cohomology $W$-modules $H^i(X_0/W)$ are free, of ranks $1,0,22,0,1$ for $i=0,1,2,3,4$.
\end{enumerate}

By [14], $H^0(X_0,T)=0$. By the definition of a K3 surface, $\Omega^2_{X_0/k}$ is trivial; hence
\begin{equation}\tag{1.1.1}
T_{X_0/k}\simeq\Omega^1_{X_0/k},
\end{equation}
and consequently $H^0(X_0,\Omega^1_{X_0/k})=0$, so that $H^2(X_0,\Omega^1_{X_0/k})=0$ by Serre duality. Moreover, $H^1(X_0,\mathcal O)=0$ and $H^2(X_0,\mathcal O)\simeq k$ by the definition of a K3 surface. Thus the table of the Hodge numbers $h^{ij}$, for $(i,j)\ne(1,1)$, is the one displayed in 1.1(a). It follows that the Hodge spectral sequence degenerates at $E_1$. Recall also (SGA~5, VII~4.11) that
\[
\sum(-1)^{i+j}h^{ij}=c_2,
\]
and that $c_2=24$ by Riemann--Roch. Hence $h^{11}=20$, which gives (b), in view of (1.1.1). It follows from (a) that the spaces $H^i_{\mathrm{DR}}(X_0/k)$ have dimensions
\newpage
\setcounter{footnote}{0}
\fontsize{10.2pt}{12.04pt}\selectfont

$1,0,22,0,1$ for $i=0,1,2,3,4$. Assertion (c) follows from the ``universal-coefficients'' exact sequence [2, VII~1.1.11]:
\[
0\longrightarrow H^i(X_0/W)\otimes k\longrightarrow H^i_{\mathrm{DR}}(X_0/k)
\longrightarrow \operatorname{Tor}^W_1(H^{i+1}(X_0/W),k)\longrightarrow0.
\]

\noindent\textbf{COROLLARY 1.2.} The versal formal space $S$ of deformations of $X_0$ over local Artinian $W$-algebras with residue field $k$ [15] is universal and formally smooth of dimension $20$; that is, it is $W$-isomorphic to $\operatorname{Spf}(W[[t_1,\ldots,t_{20}]])$.

This is an immediate consequence of 1.1(b).

In what follows in this section, we write
\begin{equation}\tag{1.3}
\UX/S
\end{equation}
for the universal deformation of $X_0$ over $S$.

\noindent\textbf{1.4.} Let $L_0$ be an invertible sheaf on $X_0$. Write $\operatorname{Def}(X_0,L_0)$ for the functor on the category $\mathcal A$ of local Artinian $W$-algebras with residue field $k$ that assigns to each $A\in\mathcal A$ the set of isomorphism classes of pairs $(X,L)$ consisting of a flat (hence smooth) deformation $X$ of $X_0$ over $A$ and an extension of $L_0$ to an invertible sheaf $L$ on $X$. Write $\operatorname{Def}(X_0)$ for the functor assigning to each $A\in\operatorname{Ob}\mathcal A$ the set of isomorphism classes of flat deformations of $X_0$ over $A$; by 1.2 this functor is (pro-)represented by $S$. There is a morphism ``forget the extension of $L_0$''
\begin{equation}\tag{1.4.1}
\operatorname{Def}(X_0,L_0)\longrightarrow\operatorname{Def}(X_0).
\end{equation}

\noindent\textbf{PROPOSITION 1.5.} The functor $\operatorname{Def}(X_0,L_0)$ is pro-representable, and the morphism (1.4.1) is a closed immersion defined by one equation.

The first assertion means that there is a largest closed formal subscheme
\begin{equation}\tag{1.5.0}
\Sigma(L_0)\subset S
\end{equation}
\newpage
\setcounter{footnote}{0}
\fontsize{10.2pt}{12.04pt}\selectfont

such that $L_0$ extends to an invertible sheaf over $\UX\times_S\Sigma(L_0)$, and such that this extension is unique.

It is readily verified that the functor $\operatorname{Def}(X_0,L_0)$ satisfies Schlessinger's conditions for the existence of a hull. Thus there is a formal scheme $S'=\operatorname{Spf}(R')$, where $R'$ is a complete Noetherian local $W$-algebra with residue field $k$, and a deformation $(\UX',\UL')$ of $(X_0,L_0)$ over $S'$ such that, for every $A\in\operatorname{Ob}\mathcal A$, the associated map
\begin{equation}\tag{1.5.1}
\operatorname{Hom}(R',A)\longrightarrow\operatorname{Def}(X_0,L_0)(A)
\end{equation}
is surjective, and is bijective for $A=k[\varepsilon]$, $\varepsilon^2=0$. Let $R$ be the ring of $S$. Since $S$ pro-represents $\operatorname{Def}(X_0)$, $\UX'$ defines a homomorphism $u:R\to R'$ such that the composite
\begin{equation}\tag{1.5.2}
\operatorname{Hom}(R',A)\xrightarrow{1.5.1}\operatorname{Def}(X_0,L_0)(A)
\xrightarrow{1.4.1}\operatorname{Def}(X_0)(A)=\operatorname{Hom}(R,A)
\end{equation}
is $\operatorname{Hom}(u,A)$. To establish the first assertion of 1.5, it is therefore enough to prove that $u$ is surjective, since then (1.5.2) is injective and hence (1.5.1) is bijective. By a familiar lemma [15, 1.1], this is equivalent to proving that, if $\mathfrak m$ (respectively $\mathfrak m'$) is the maximal ideal of $R$ (respectively $R'$), the map
$\mathfrak m/(pR+\mathfrak m^2)\to\mathfrak m'/(pR'+\mathfrak m'^2)$ induced by $u$ is surjective; equivalently, that the tangent map at the origin of (1.4.1),
\begin{equation}\tag{1.5.3}
\operatorname{Hom}(u,k[\varepsilon]):\operatorname{Def}(X_0,L_0)(k[\varepsilon])
\longrightarrow\operatorname{Def}(X_0)(k[\varepsilon]),
\end{equation}
is injective. Let $T'$ be the sheaf on $X_0$ of automorphisms of the trivial deformation of $(X_0,L_0)$ over $k[\varepsilon]$. There is an exact sequence
\begin{equation}\tag{1.5.4}
0\longrightarrow\mathcal O_{X_0}\longrightarrow T'\longrightarrow T\longrightarrow0,
\end{equation}
where $T'\to T$ is defined by forgetting $L_0$, and $T=T_{X_0}$ is viewed as the sheaf of automorphisms of the trivial deformation of $X_0$ over $k[\varepsilon]$. It is standard that (1.5.3) identifies canonically with the map
\newpage
\setcounter{footnote}{0}
\fontsize{10.2pt}{12.04pt}\selectfont

$H^1(X_0,T')\to H^1(X_0,T)$ induced by $T'\to T$. Since $H^1(X_0,\mathcal O)=0$, the long exact cohomology sequence of (1.5.4) shows that (1.5.3) is injective. It remains to prove that the closed immersion $S'\hookrightarrow S$ is defined by a single equation, i.e. that $I=\ker(u)$ is principal. Put $S''=\operatorname{Spf}(R/\mathfrak m I)$; this is a square-zero thickening of $S'$ in $S$, with ideal $I/\mathfrak m I$. The obstruction to extending the invertible sheaf $\UL'$ above $\UX\times_S S''$ is an element
\[
a\in H^2(X_0,I/\mathfrak m I)=H^2(X_0,\mathcal O)\otimes I/\mathfrak m I,
\]
which, after choosing a basis of $H^2(X_0,\mathcal O)$, we regard as an element of $I/\mathfrak m I$. Let $\Sigma=\operatorname{Spf}(R/(\mathfrak m I+(f)))$, where $f\in I$ lifts $a$. Then $S'\subset\Sigma\subset S''\subset S$, and by construction (and functoriality of the obstruction) $\UL'$ extends to $\UX\times_S\Sigma$. The universal property of $S'$ therefore implies $S'=\Sigma$, i.e. $\mathfrak m I+(f)=I$. Nakayama's lemma now shows that $f$ generates $I$, completing the proof.

We may now state the main result of the exposé.

\noindent\textbf{THEOREM 1.6.} Let $L_0$ be a nontrivial invertible sheaf on $X_0$. With the notation of 1.5, the formal scheme $\Sigma(L_0)$ pro-representing $\operatorname{Def}(X_0,L_0)$ is flat over $W$, of relative dimension $19$.

In other words, if $f$ is an equation for $\Sigma(L_0)$ in $S$ (1.5), then $p$ does not divide $f$. Equivalently, $\Sigma(L_0)$ does not contain the reduction $S_0$ of $S$ modulo $p$, i.e. $L_0$ cannot extend over $\UX\times_S S_0$.

The proof of 1.6 will be given in Section~2. We end this section with a few consequences of 1.6.

\noindent\textbf{COROLLARY 1.7.} Let $L_0$ be a nontrivial invertible sheaf on $X_0$. There exist a trait $T$ finite over $W$, a deformation of $X_0$ to a formal scheme $\UX$ flat over $T$, and an extension of $L_0$ to an invertible sheaf $\UL$ on $\UX$.
\newpage
\setcounter{footnote}{0}
\fontsize{10.2pt}{12.04pt}\selectfont

We must prove that there is a $W$-morphism $T\to\Sigma(L_0)$, where $T$ is a trait finite over $W$. Since $p$ is not a zero divisor in the ring $R'$ of $\Sigma(L_0)$, there exist (EGA $0_{\mathrm{IV}}$ 16.4.1) elements $x_1,\ldots,x_n$ of the maximal ideal of $R'$ which, together with $p$, form a system of parameters for $R'$ (thus $n=19$). The quotient $B=R'/(x_1,\ldots,x_n)$ is quasi-finite over $W$, hence finite over $W$. There is therefore a local $W$-homomorphism from $B$ to a complete discrete valuation ring $C$ finite over $W$, and the composite $R'\to B\to C$ gives the required morphism.

Applying Grothendieck's algebraization theorem (EGA~III~5.4.5), one deduces from 1.7 the announced lifting theorem.

\noindent\textbf{COROLLARY 1.8.} Let $L_0$ be an ample invertible sheaf on $X_0$. There exist a trait $T$ finite over $W$, a deformation of $X_0$ to a proper smooth scheme $X$ over $T$, and an extension of $L_0$ to an ample invertible sheaf $L$ on $X$.

\noindent\textbf{REMARK 1.9.} It is not known whether every K3 surface over $k$ lifts to a proper smooth scheme over $W$. Ogus [13, \S2] shows that: (a) every K3 surface over $k$ lifts over $W$, except possibly in the ``superspecial'', non-elliptic case, which in fact does not exist if one assumes Artin's conjecture [1]; (b) if $p>2$, every K3 surface over $k$ lifts over $W[\sqrt p]$. Thus, at present, only the special case of 1.8 in which $p=2$ and $X_0$ is superspecial is not covered by these results. We also point out that the cited paper of Ogus contains useful further information on the structure of $\Sigma(L_0)$; see also the following exposé for the case in which $X_0$ is ordinary.

\noindent\textbf{COROLLARY 1.10.} If $k$ is the algebraic closure of a finite field over which $X_0$ is defined, the corresponding Frobenius acts semisimply on $H^2(X_0,\mathbf Q_\ell)$, where $\ell\ne p$.
\newpage
\setcounter{footnote}{0}
\fontsize{10.2pt}{12.04pt}\selectfont

This follows from [5]. With the notation of 1.8, the generic fibre of $X$ is a K3 surface (the property of a surface of being K3 is stable under generization, since it is expressed by ``$K$ algebraically equivalent to zero and $\chi(\mathcal O)=2$''). Hence $X_0$ satisfies the hypothesis of [5, 1.2]; the conclusion follows from [5, 6.6] and from the fact that Frobenius acts semisimply on the $\ell$-adic $H^1$ of an abelian variety over a finite field ([16], [11, p.~203]).

Notice that 1.10 is in fact independent of the assertion $H^0(X_0,T_{X_0})=0$: if $X_0$ admitted a nonzero vector field, then $X_0$ would be unirational by [14], and the conclusion of 1.10 would still hold (trace argument).

\section*{2. DE RHAM COHOMOLOGY OF $\UX/S$ AND PROOF OF 1.6}
We retain the notation of Section~1: $X_0$ is a K3 surface over $k$, and $\UX/S$ denotes its universal $W$-deformation (1.3). Readers familiar with de Rham cohomology may skip 2.1--2.10, which merely recall standard facts concerning the Gauss--Manin connection, the Hodge filtration, the Frobenius operation, and the Chern class of an invertible sheaf.

\noindent\textbf{2.1.} Let $\Omega^\bullet_{\UX/S}$ be the de Rham complex of the relative formal scheme $\UX/S$ (by definition, for each $i$, $\Omega^i_{\UX/S}$ is the inverse limit of the usual modules of differentials $\Omega^i_{\UX\times_S S'/S'}$, as $S'$ runs through the infinitesimal neighbourhoods of $\operatorname{Spec}(k)$ in $S$). Let $f:\UX\to S$ be the projection. By definition, the de Rham cohomology of $\UX/S$ consists of the $\mathcal O_S$-modules
\begin{equation}\tag{2.1.1}
H^i_{\mathrm{DR}}(\UX/S)\overset{\mathrm{def}}{=}R^if_*(\Omega^\bullet_{\UX/S}),
\end{equation}
whereas the Hodge cohomology of $\UX/S$ consists of the $\mathcal O_S$-modules
\begin{equation}\tag{2.1.2}
H^i(\UX,\Omega^j_{\UX/S})\overset{\mathrm{def}}{=}R^if_*(\Omega^j_{\UX/S}).
\end{equation}
\newpage
\setcounter{footnote}{0}
\fontsize{10.2pt}{12.04pt}\selectfont

Since the $\mathcal O_{\UX}$-modules $\Omega^i_{\UX/S}$ are locally free of finite type, the $\mathcal O_S$-modules (2.1.1) and (2.1.2) are coherent by Grothendieck's finiteness theorem (EGA~III~3.4.2). There is also the usual spectral sequence (the ``Hodge spectral sequence'')
\begin{equation}\tag{2.1.3}
E_1^{ij}=H^j(\UX,\Omega^i_{\UX/S})\Longrightarrow H^*_{\mathrm{DR}}(\UX/S).
\end{equation}

\noindent\textbf{PROPOSITION 2.2.}
\begin{enumerate}[label=\alph*)]
\item The spectral sequence (2.1.3) degenerates at $E_1$; the $\mathcal O_S$-modules $H^j(\UX,\Omega^i_{\UX/S})$ are free of finite type, and the canonical maps
\[
H^j(\UX,\Omega^i_{\UX/S})\otimes k\longrightarrow H^j(X_0,\Omega^i_{X_0/k})
\]
are isomorphisms.
\item The $\mathcal O_S$-modules $H^i_{\mathrm{DR}}(\UX/S)$ are free of finite type, and the canonical maps
\[
H^i_{\mathrm{DR}}(\UX/S)\otimes k\longrightarrow H^i_{\mathrm{DR}}(X_0/k)
\]
are isomorphisms.
\item The cup product
\[
\langle\ ,\ \rangle:H^2_{\mathrm{DR}}(\UX/S)\otimes H^2_{\mathrm{DR}}(\UX/S)
\longrightarrow H^4_{\mathrm{DR}}(\UX/S)
\]
is a perfect duality.
\end{enumerate}

Because the table of Hodge numbers of $X_0/k$ is ``interlaced with zeros'' (1.1(a)), the usual criterion (EGA~III~7.5.4), applied to the cohomological functors
$M\mapsto H^\bullet(\UX,\Omega^i_{\UX/S}\otimes f^*M)$, immediately gives the second assertion of (a), and the first follows. Assertion (b) follows from (a), and implies (c) by Poincaré duality for $H^*_{\mathrm{DR}}(X_0/k)$.

\noindent\textbf{2.3.} Let $\Omega^\bullet_{S/W}$ be the de Rham complex of the ``formal'' differentials of $S/W$: $\Omega^i_{S/W}=\bigwedge^i\Omega^1_{S/W}$, where $\Omega^1_{S/W}$, the inverse limit of the completed modules of differentials $\Omega^1_{S_n/W_n}$ (with $S_n/W_n$ the reduction of $S/W$ modulo $p^{n+1}$), is free over $\mathcal O_S$ with basis $dt_1,\ldots,dt_{20}$ if $S\simeq W[[t_1,\ldots,t_{20}]]$. The $H^i_{\mathrm{DR}}(\UX/S)$ carry a canonical integrable connection, the Gauss--Manin connection,
\begin{equation}\tag{2.3.1}
\nabla:H^i_{\mathrm{DR}}(\UX/S)\longrightarrow\Omega^1_{S/W}\otimes H^i_{\mathrm{DR}}(\UX/S).
\end{equation}
The simplest definition of (2.3.1) paraphrases the Katz--Oda construction [\,], starting from the canonical extension
\newpage
\setcounter{footnote}{0}
\fontsize{10.2pt}{12.04pt}\selectfont

of $\mathcal O_{\UX}$-modules
\begin{equation}\tag{2.3.0}
0\longrightarrow f^*\Omega^1_{S/W}\longrightarrow\Omega^1_{\UX/W}\longrightarrow\Omega^1_{\UX/S}\longrightarrow0,
\end{equation}
where $\Omega^1_{\UX/W}$ is the module of completed differentials. One may also use the fact that $H^i_{\mathrm{DR}}(\UX/S)$ is the value at $S$ of a crystal in $\mathcal O$-modules on the crystalline site of $S_0/W$:
\begin{equation}\tag{2.3.2}
H^i_{\mathrm{DR}}(\UX/S)=R^i(f_0)_{\mathrm{cris}*}(\mathcal O_{\UX_0/W})(S),
\end{equation}
where
\begin{equation}\tag{2.3.2.1}
(f_0:\UX_0\to S_0)=f\times_S S_0,\qquad
S_0=S\otimes_W k\simeq\operatorname{Spf}(k[[t_1,\ldots,t_{20}]]),
\end{equation}
and $(f_0)_{\mathrm{cris}}:(\UX_0/W)_{\mathrm{cris}}\to(S_0/W)_{\mathrm{cris}}$ is the corresponding morphism of crystalline sites. This is a ``formal variant'' of Berthelot's result [2, V~3.6], readily deduced from it by applying it to the morphisms induced by $f_0$ over the infinitesimal neighbourhoods of $\operatorname{Spec}(k)$ in $S_0$.

Either definition shows that the cup product of 2.2(c) is horizontal:
\begin{equation}\tag{2.3.3}
\langle\nabla x,y\rangle+\langle x,\nabla y\rangle=\nabla\langle x,y\rangle
\end{equation}
for all $x,y\in H^2_{\mathrm{DR}}(\UX/S)$.

Let $F^i_{\mathrm{Hdg}}H^2_{\mathrm{DR}}(\UX/S)$ be the Hodge filtration, the abutment of (2.1.3). If $\Omega^{\ge i}_{\UX/S}=(0\to\Omega^i_{\UX/S}\to\Omega^{i+1}_{\UX/S}\to\cdots)$ denotes the de Rham complex truncated by zero in degrees $<i$, then the inclusion $\Omega^{\ge i}_{\UX/S}\to\Omega^\bullet_{\UX/S}$ gives, by 2.2(a), an isomorphism
\begin{equation}\tag{2.3.4}
H^2(\UX,\Omega^{\ge i}_{\UX/S})\simeq F^i_{\mathrm{Hdg}}H^2_{\mathrm{DR}}(\UX/S).
\end{equation}
We have
\begin{equation}\tag{2.3.5}
H^2_{\mathrm{DR}}(\UX/S)=F^0_{\mathrm{Hdg}}\supset F^1_{\mathrm{Hdg}}
\supset F^2_{\mathrm{Hdg}}\supset F^3_{\mathrm{Hdg}}=0,
\end{equation}
the $F^i_{\mathrm{Hdg}}$ are free $\mathcal O_S$-modules, and degeneration of (2.1.3) gives isomorphisms
\begin{equation}\tag{2.3.6}
H^{2-i}(\UX,\Omega^i_{\UX/S})\simeq\operatorname{gr}^i_F H^2_{\mathrm{DR}}(\UX/S),
\qquad(F=F_{\mathrm{Hdg}}).
\end{equation}
\newpage
\setcounter{footnote}{0}
\fontsize{10.2pt}{12.04pt}\selectfont

In particular, $F^1_{\mathrm{Hdg}}$ (respectively $F^2_{\mathrm{Hdg}}$) is free of rank $21$ (respectively $1$). From (2.3.4) one sees that the orthogonal of $F^2_{\mathrm{Hdg}}$ for the cup product of 2.2(c) contains $F^1_{\mathrm{Hdg}}$, and therefore equals it:
\begin{equation}\tag{2.3.7}
F^1_{\mathrm{Hdg}}=(F^2_{\mathrm{Hdg}})^\perp
\qquad\text{(and hence }F^2_{\mathrm{Hdg}}=(F^1_{\mathrm{Hdg}})^\perp\text{).}
\end{equation}
The Katz--Oda description of the Gauss--Manin connection shows that
\begin{equation}\tag{2.3.8}
\nabla F^i_{\mathrm{Hdg}}H^2_{\mathrm{DR}}(\UX/S)
\subset\Omega^1_{S/W}\otimes F^{i-1}_{\mathrm{Hdg}}H^2_{\mathrm{DR}}(\UX/S)
\end{equation}
(``Griffiths transversality''; see, for example, [8, 1.4]). Passing to quotients, $\nabla$ therefore induces an $\mathcal O_S$-linear map
\begin{equation}\tag{2.3.9}
\operatorname{gr}^i\nabla:\operatorname{gr}^i_FH^2_{\mathrm{DR}}(\UX/S)
\longrightarrow\Omega^1_{S/W}\otimes\operatorname{gr}^{i-1}_F H^2_{\mathrm{DR}}(\UX/S).
\end{equation}
Correspondingly there is an $\mathcal O_S$-linear map
\begin{equation}\tag{2.3.10}
\nabla_i:T_{S/W}\longrightarrow
\operatorname{Hom}(\operatorname{gr}^i_FH^2_{\mathrm{DR}}(\UX/S),
\operatorname{gr}^{i-1}_FH^2_{\mathrm{DR}}(\UX/S)),
\end{equation}
where $T_{S/W}=(\Omega^1_{S/W})^\vee$. By (2.3.6), the target of (2.3.10) identifies canonically, for the fixed $i$, with
\[
\operatorname{Hom}\bigl(H^{2-i}(\UX,\Omega^i_{\UX/S}),
H^{3-i}(\UX,\Omega^{i-1}_{\UX/S})\bigr).
\]
Let $T_{\UX/S}=(\Omega^1_{\UX/S})^\vee$. If
\begin{equation}\tag{2.3.11}
\operatorname{Kod}(\UX/S):T_{S/W}\longrightarrow H^1(\UX,T_{\UX/S})
\end{equation}
denotes the Kodaira--Spencer map associated with $\UX/S/W$, defined by the extension (2.3.0) (cf.~[8, 1.3]), then, under the preceding identification, (2.3.10) belongs to the commutative triangle
\begin{equation}\tag{2.3.12}
\begin{tikzcd}[column sep=large,row sep=large]
T_{S/W}\arrow[r,"\operatorname{Kod}(\UX/S)"]\arrow[dr,"\nabla_i"']&H^1(\UX,T_{\UX/S})\arrow[d]\\
&\operatorname{Hom}(H^{2-i}(\UX,\Omega^i_{\UX/S}),H^{3-i}(\UX,\Omega^{i-1}_{\UX/S})),
\end{tikzcd}
\end{equation}
where the vertical arrow is defined by cup product, via the interior product $T_{\UX/S}\otimes\Omega^i_{\UX/S}\to\Omega^{i-1}_{\UX/S}$. This compatibility is verified exactly as
\newpage
\setcounter{footnote}{0}
\fontsize{10.2pt}{12.04pt}\selectfont

in [8, 1.4.1.7].

By (2.3.7), the cup product of 2.2(c), after passage to $\operatorname{gr}_F$, induces a perfect duality
\begin{equation}\tag{2.3.13}
\operatorname{gr}^i_FH^2_{\mathrm{DR}}(\UX/S)\otimes
\operatorname{gr}^{2-i}_FH^2_{\mathrm{DR}}(\UX/S)
\longrightarrow H^4_{\mathrm{DR}}(\UX/S)\simeq H^2(\UX,\Omega^2_{\UX/S}),
\end{equation}
corresponding under (2.3.6) to the cup product
\[
H^{2-i}(\UX,\Omega^i_{\UX/S})\otimes H^i(\UX,\Omega^{2-i}_{\UX/S})
\longrightarrow H^2(\UX,\Omega^2_{\UX/S}).
\]
It follows from (2.3.3) that, for every $D\in T_{S/W}$,
\begin{equation}\tag{2.3.14}
\nabla_1(D)=-\nabla_2(D)^\vee,
\end{equation}
i.e. $\langle\nabla_2(D)x,y\rangle+\langle x,\nabla_1(D)y\rangle=0$ for $x\in\operatorname{gr}^2_F$, $y\in\operatorname{gr}^1_F$.

The following statement expresses the mobility of the Hodge filtration under Gauss--Manin.

\noindent\textbf{PROPOSITION 2.4.} The maps (2.3.10) and (2.3.11) are isomorphisms.

First, the sheaf $\Omega^2_{\UX/S}$, the extension to $\UX$ of the trivial sheaf $\Omega^2_{X_0/k}$, is necessarily trivial (1.5); hence
\begin{equation}\tag{2.4.1}
T_{\UX/S}\simeq\Omega^1_{\UX/S}.
\end{equation}
Consequently, by 2.2(a), $H^1(\UX,T_{\UX/S})$ is free of rank $20$, and the canonical map $H^1(\UX,T_{\UX/S})\otimes k\to H^1(X_0,T_{X_0/k})$ is an isomorphism. A standard argument shows that
\begin{equation}\tag{2.4.2}
\operatorname{Kod}(\UX/S)\otimes_Wk:T_{S/W}\otimes_Wk\longrightarrow H^1(X_0,T_{X_0/k})
\end{equation}
is precisely the tangent map at the origin of the canonical isomorphism $S\simeq\operatorname{Def}(X_0)$ (1.4), so $\operatorname{Kod}(\UX/S)$ is an isomorphism. By (2.3.12), it remains to prove that the vertical arrow in (2.3.12) is an isomorphism; by (2.3.14), it is enough to do so for $i=1$. By 2.2(a) and the observation at the beginning of the proof, it suffices to show that the cup product
\newpage
\setcounter{footnote}{0}
\fontsize{10.2pt}{12.04pt}\selectfont

\begin{equation}\tag{2.4.3}
H^1(X_0,T_{X_0/k})\otimes H^1(X_0,\Omega^1_{X_0/k})
\longrightarrow H^2(X_0,\mathcal O)
\end{equation}
is nondegenerate. But the choice of a basis of $H^0(X_0,\Omega^2_{X_0/k})$ identifies $T_{X_0/k}$ with $\Omega^1_{X_0/k}$ and identifies (2.4.3) with Serre duality
\[
H^1(X_0,\Omega^1_{X_0/k})\otimes H^1(X_0,\Omega^1_{X_0/k})
\longrightarrow H^2(X_0,\Omega^2_{X_0/k}),
\]
which proves the assertion.

\noindent\textbf{COROLLARY 2.4.4.} The map $\operatorname{gr}^1\nabla$ (respectively $\operatorname{gr}^2\nabla$) in (2.3.9) is an isomorphism (respectively is injective with free cokernel).

Consequently the same assertion holds for $\operatorname{gr}^1(\nabla|S_0)$ (respectively $\operatorname{gr}^2(\nabla|S_0)$), where $\nabla|S_0$ is the Gauss--Manin connection on $H^2_{\mathrm{DR}}(\UX_0/S_0)$. In view of the duality (2.3.13), 2.4.4 follows immediately from the fact that (2.3.10) is an isomorphism.

\noindent\textbf{2.5.} With the notation of (2.3.2.1), let $F_{\UX_0/S_0}:\UX_0\to \UX_0^{(p)}$ be the relative Frobenius of $\UX_0/S_0$, defined by the usual Cartesian diagram
\begin{equation}\tag{2.5.1}
\begin{tikzcd}[column sep=large,row sep=2em]
\UX_0\arrow[d] & \UX_0^{(p)}\arrow[l]\arrow[d] & \UX_0\arrow[l,"F_{\UX_0/S_0}"']\arrow[dl]\\
S_0 & S_0\arrow[l] &
\end{tikzcd}
\end{equation}
where the horizontal composites are the absolute Frobenius maps. The crystalline interpretation (2.3.2) of $H^i_{\mathrm{DR}}(\UX/S)$ shows, by functoriality of crystalline cohomology, that $F_{\UX_0/S_0}$ induces, for every lift $\varphi:S\to S$ of the absolute Frobenius of $S_0$ compatible with the canonical Frobenius $\sigma$ of $W$, an $\mathcal O_S$-linear horizontal homomorphism
\begin{equation}\tag{2.5.2}
F(\varphi):\varphi^*H^i_{\mathrm{DR}}(\UX/S)\longrightarrow H^i_{\mathrm{DR}}(\UX/S).
\end{equation}
\newpage
\setcounter{footnote}{0}
\fontsize{10.2pt}{12.04pt}\selectfont

This homomorphism is an isogeny (i.e. $F(\varphi)\otimes_{\mathbf Z_p}\mathbf Q_p$ is an isomorphism), and its dependence on $\varphi$ is controlled by $\nabla$ (cf.~[7] and the following exposé). In what follows we fix a lift $\varphi$, for example the one given by $\varphi(t_i)=t_i^p$ ($1\le i\le20$), with $S=\operatorname{Spf}W[[t_1,\ldots,t_{20}]]$, so that
$\varphi(\sum a_{\alpha}t^{\alpha})=\sum a_{\alpha}^{\sigma}t^{p\alpha}$. We write $F(\varphi)=\widetilde F$ and denote by
\begin{equation}\tag{2.5.3}
F:H^i_{\mathrm{DR}}(\UX/S)\longrightarrow H^i_{\mathrm{DR}}(\UX/S)
\end{equation}
the $\sigma$-linear homomorphism obtained by composing $\widetilde F$ with the adjunction map
$H^i_{\mathrm{DR}}(\UX/S)\to\varphi^*H^i_{\mathrm{DR}}(\UX/S)$, $x\mapsto1\otimes x$. By construction, $F$ is compatible with the cup product of 2.2(c):
\begin{equation}\tag{2.5.4}
\langle Fx,Fy\rangle=F\langle x,y\rangle
\end{equation}
for all $x,y\in H^2_{\mathrm{DR}}(\UX/S)$.

The following statement is a special case of a theorem of Mazur--Ogus ([10], [4, 8.26])\footnote{This is a formal variant of [4, 8.26], from which it follows easily.}.

\noindent\textbf{PROPOSITION 2.6.} With the notation of 2.3 and 2.5,
\[
F^1_{\mathrm{Hdg}}H^2_{\mathrm{DR}}(\UX/S)
\subset\{x\in H^2_{\mathrm{DR}}(\UX/S)\mid Fx\in pH^2_{\mathrm{DR}}(\UX/S)\},
\]
and the two sides have the same image in $H^2_{\mathrm{DR}}(\UX_0/S_0)$.

In other words, if $F:H^2_{\mathrm{DR}}(\UX_0/S_0)\to H^2_{\mathrm{DR}}(\UX_0/S_0)$ denotes the $p$-linear Frobenius endomorphism (the reduction modulo $p$ of (2.5.3)), and if $F^i_{\mathrm{Hdg}}H^2_{\mathrm{DR}}(\UX_0/S_0)$ is the Hodge filtration (the reduction modulo $p$ of that on $H^2_{\mathrm{DR}}(\UX/S)$), then
\begin{equation}\tag{2.6.1}
F^1_{\mathrm{Hdg}}H^2_{\mathrm{DR}}(\UX_0/S_0)
=\ker\bigl(F:H^2_{\mathrm{DR}}(\UX_0/S_0)\to H^2_{\mathrm{DR}}(\UX_0/S_0)\bigr).
\end{equation}
We recall the proof. By definition, $F$ is the composite of the canonical $p$-linear injection
$H^2_{\mathrm{DR}}(\UX_0/S_0)\to H^2_{\mathrm{DR}}(\UX_0^{(p)}/S_0)=F_{S_0}^*H^2_{\mathrm{DR}}(\UX_0/S_0)$ defined by the Cartesian square (2.5.1), and the $\mathcal O_{S_0}$-linear map
$\widetilde F:H^2_{\mathrm{DR}}(\UX_0^{(p)}/S_0)\to H^2_{\mathrm{DR}}(\UX_0/S_0)$ defined by the relative Frobenius $F_{\UX_0/S_0}$. It is immediate that
\newpage
\setcounter{footnote}{0}
\fontsize{10.2pt}{12.04pt}\selectfont

\[
F^1_{\mathrm{Hdg}}H^2_{\mathrm{DR}}(\UX_0/S_0)=
H^2_{\mathrm{DR}}(\UX_0/S_0)\cap F^1_{\mathrm{Hdg}}H^2_{\mathrm{DR}}(\UX_0^{(p)}/S_0),
\]
so it is enough to prove
\begin{equation}\tag{2.6.2}
F^1_{\mathrm{Hdg}}H^2_{\mathrm{DR}}(\UX_0^{(p)}/S_0)
=\ker\bigl(\widetilde F:H^2_{\mathrm{DR}}(\UX_0^{(p)}/S_0)\to H^2_{\mathrm{DR}}(\UX_0/S_0)\bigr).
\end{equation}
By 2.2, the Hodge spectral sequence of $\UX_0^{(p)}/S_0$ degenerates at $E_1$, and therefore there is an exact sequence
\begin{equation}\tag{2.6.3}
0\longrightarrow F^1_{\mathrm{Hdg}}H^2_{\mathrm{DR}}(\UX_0^{(p)}/S_0)
\longrightarrow H^2_{\mathrm{DR}}(\UX_0^{(p)}/S_0)
\xrightarrow{\pi}H^2(\UX_0^{(p)},\mathcal O)\longrightarrow0,
\end{equation}
where $\pi$ is the canonical projection. Since $F_{\UX_0/S_0}$ vanishes on $\Omega^i_{\UX_0^{(p)}/S_0}$ for $i\ge1$, there is a commutative square
\[
\begin{tikzcd}[column sep=large,row sep=large]
H^2_{\mathrm{DR}}(\UX_0^{(p)}/S_0)\arrow[r,"\pi"]\arrow[d,"\widetilde F"']&H^2(\UX_0^{(p)},\mathcal O)\arrow[d,"\widetilde F"]\\
H^2_{\mathrm{DR}}(\UX_0/S_0)&H^2(\UX_0,\mathcal H^0(\Omega^\bullet_{\UX_0/S_0}))\arrow[l,"(*)"']
\end{tikzcd}
\]
where $(*)$ is induced by the inclusion $\mathcal H^0(\Omega^\bullet_{\UX_0/S_0})\hookrightarrow\Omega^\bullet_{\UX_0/S_0}$. Degeneration at $E_1$ of the Hodge spectral sequence of $\UX_0/S_0$ entails, by the standard argument using the Cartier operator, degeneration at $E_2$ of the conjugate spectral sequence
\[
E_2^{ij}=H^i(\UX_0,\mathcal H^j(\Omega^\bullet_{\UX_0/S_0}))\Longrightarrow H^*_{\mathrm{DR}}(\UX_0/S_0).
\]
In particular, $(*)$ is injective. The right vertical arrow is an isomorphism (a special case of the Cartier isomorphism), so $\widetilde F$ and $\pi$ have the same kernel. Together with (2.6.3), this proves the assertion.

\noindent\textbf{REMARKS 2.7.}

\noindent a) By [4, 8.26], one also has
\begin{equation}\tag{2.7.1}
F^2_{\mathrm{Hdg}}H^2_{\mathrm{DR}}(\UX_0/S_0)
=\operatorname{Im}\!\left(
\{x\in H^2_{\mathrm{DR}}(\UX/S)\mid Fx\in p^2H^2_{\mathrm{DR}}(\UX/S)\}
\longrightarrow H^2_{\mathrm{DR}}(\UX_0/S_0)\right),
\end{equation}
but we shall not need this.
\noindent b) The $F$-crystal $H^4_{\mathrm{DR}}(\UX/S)$ is isomorphic to the ``Tate crystal'' $\mathcal O(-2)$, namely $\mathcal O_S$ equipped with $\nabla=d$ and $F=p^2\varphi$. In the absence, in the literature, of a trace morphism $H^4_{\mathrm{DR}}(\UX/S)\to\mathcal O_S$, one may nevertheless
\newpage
\setcounter{footnote}{0}
\fontsize{10.2pt}{12.04pt}\selectfont

see this as follows. First, $H^4_{\mathrm{DR}}(\UX/S)$ is the image under cup product of
$F^1H^2_{\mathrm{DR}}(\UX/S)\otimes F^1H^2_{\mathrm{DR}}(\UX/S)$. By (2.5.4) and 2.6, it follows that
$F(H^4_{\mathrm{DR}}(\UX/S))\subset p^2H^4_{\mathrm{DR}}(\UX/S)$. Since
$H^4_{\mathrm{DR}}(\UX/S)\otimes W=H^4(X_0/W)$ is isomorphic to $W(-2)$, the crystal $H^4_{\mathrm{DR}}(\UX/S)$ has the form $U(-2)$, where $U$ is a unit-root $F$-crystal; every unit-root $F$-crystal over $S$ is trivial [7], proving the assertion.

\noindent\textbf{2.8.} The last ingredient needed in the proof of 1.6 is the crystalline Chern class of an invertible sheaf (cf.~[3]). Let $\operatorname{Fil}^1\Omega^\bullet_{\UX/S}$ be the subcomplex of $\Omega^\bullet_{\UX/S}$ which is the kernel of the canonical projection $\Omega^\bullet_{\UX/S}\to\mathcal O_{\UX_0}$, i.e.
\begin{equation}\tag{2.8.1}
\operatorname{Fil}^1\Omega^\bullet_{\UX/S}
=(p\mathcal O_{\UX}\xrightarrow{d}\Omega^1_{\UX/S}\xrightarrow{d}\Omega^2_{\UX/S}).
\end{equation}
The defining exact sequence for $\operatorname{Fil}^1$,
\begin{equation}\tag{2.8.2}
0\longrightarrow\operatorname{Fil}^1\Omega^\bullet_{\UX/S}
\longrightarrow\Omega^\bullet_{\UX/S}\longrightarrow\mathcal O_{\UX_0}\longrightarrow0,
\end{equation}
gives an exact sequence
\begin{equation}\tag{2.8.3}
0\longrightarrow1+\operatorname{Fil}^1\Omega^\bullet_{\UX/S}
\longrightarrow(\Omega^\bullet_{\UX/S})^*\longrightarrow\mathcal O_{\UX_0}^*\longrightarrow0,
\end{equation}
where
\[
(\Omega^\bullet_{\UX/S})^*=(\mathcal O_{\UX}^*\xrightarrow{d\log}\Omega^1_{\UX/S}\xrightarrow{d}\Omega^2_{\UX/S})
\]
is the ``multiplicative'' de Rham complex. By definition, the Chern-class homomorphism
\begin{equation}\tag{2.8.4}
c_1:\operatorname{Pic}(\UX_0)=H^1(\UX_0,\mathcal O_{\UX_0}^*)
\longrightarrow H^2(\UX,\operatorname{Fil}^1\Omega^\bullet_{\UX/S})
\end{equation}
is the composite of the connecting morphism of (2.8.3) and the map induced by the morphism of complexes
\begin{equation}\tag{2.8.5}
\log:1+\operatorname{Fil}^1\Omega^\bullet_{\UX/S}
\longrightarrow\operatorname{Fil}^1\Omega^\bullet_{\UX/S},
\end{equation}
which is the identity in degrees $\ge1$ and sends
$1+x$ to $\sum(-1)^ix^{i+1}/(i+1)$ in degree zero. The legitimacy of (2.8.5) is immediate.
\newpage
\setcounter{footnote}{0}
\fontsize{10.2pt}{12.04pt}\selectfont

Since $H^1(\UX_0,\mathcal O)=0$ (2.2(a)), the canonical map
\begin{equation}\tag{2.8.6}
H^2(\UX,\operatorname{Fil}^1\Omega^\bullet_{\UX/S})
\longrightarrow H^2_{\mathrm{DR}}(\UX/S)
\end{equation}
is injective. We shall therefore sometimes regard $c_1$ in (2.8.4) as taking values in $H^2_{\mathrm{DR}}(\UX/S)$.

\noindent\textbf{PROPOSITION 2.9.} Let $\UL_0\in\operatorname{Pic}(\UX_0)$ and put $x=c_1(\UL_0)$. Then
\begin{enumerate}[label=\alph*)]
\item $Fx=px$;
\item $\nabla x=0$.
\end{enumerate}

For these formulas it is convenient to interpret $c_1$ crystallinely. The construction of [3, 2.1] gives a homomorphism
\begin{equation}\tag{2.9.1}
c_{1/W}:\operatorname{Pic}(\UX_0)\longrightarrow H^2_{\mathrm{cris}}(\UX_0,J_{\UX_0/W}),
\end{equation}
where $J_{\UX_0/W}$ is the crystalline sheaf defined by the canonical exact sequence
\begin{equation}\tag{2.9.2}
0\longrightarrow J_{\UX_0/W}\longrightarrow\mathcal O_{\UX_0/W}\longrightarrow\mathcal O_{\UX_0}\longrightarrow0.
\end{equation}
By (2.3.2),
\[
H^2(\UX,\operatorname{Fil}^1\Omega^\bullet_{\UX/S})
=R^2(f_0)_{\mathrm{cris}*}(J_{\UX_0/W})(S),
\]
and hence
\begin{equation}\tag{2.9.3}
H^0_{\mathrm{cris}}(S_0,R^2(f_0)_{\mathrm{cris}*}J_{\UX_0/W})
=\{x\in H^2(\UX,\operatorname{Fil}^1\Omega^\bullet_{\UX/S})\mid\nabla x=0\}.
\end{equation}
It is readily checked that (2.8.4) is the composite of (2.9.1) and the canonical homomorphism
\begin{equation}\tag{2.9.4}
H^2_{\mathrm{cris}}(\UX_0,J_{\UX_0/W})\longrightarrow
H^0_{\mathrm{cris}}(S_0,R^2(f_0)_{\mathrm{cris}*}J_{\UX_0/W}),
\end{equation}
under the identification (2.9.3). This proves 2.9(b). The preceding description also shows functoriality in $\UX_0/S_0$ of (2.8.4), which implies 2.9(a), because, for the absolute Frobenius $F_{\UX_0}$ of $\UX_0$,
\[
Fx=c_1(F_{\UX_0}^*\UL_0)=c_1(\UL_0^{\otimes p})=px.
\]
\newpage
\setcounter{footnote}{0}
\fontsize{10.2pt}{12.04pt}\selectfont

\noindent\textbf{REMARK 2.9.5.} As regards (b), a reader reluctant to use crystalline considerations may verify it ``by hand'' from the explicit Katz--Oda description [9] of the Gauss--Manin connection. Let $D$ be a derivation of $S/W$. Choose an affine open covering $\mathcal U$ of $\UX_0$, a cocycle $(g_{ij})$ with values in $\mathcal O_{\UX_0}^*$ representing $\UL_0$, and lifts $D_i$ of $D$ to derivations of $\UX/W$ on $U_i$. Lifting $g_{ij}$ to a section $\widetilde g_{ij}$ of $\mathcal O_{\UX}^*$ on $U_i\cap U_j$ defines a $2$-cocycle $h_{ijk}$ with values in $1+p\mathcal O_{\UX}$; the class $x=c_1(\UL_0)$ is represented by
\[
(\log h_{ijk})+(d\log g_{ij})\in Z^2(\mathcal U,\operatorname{Fil}^1\Omega^\bullet_{\UX/S}).
\]
On the other hand, by [9], $\nabla(D)x$ is the class of the $2$-cocycle ($i<j<k$)
\[
y=(D_i\log h_{ijk})+(dD_i\log g_{ij})-((D_i\log-D_j\log)g_{jk}).
\]
One checks that $y$ is the coboundary of the $1$-cochain
$(D_i\log\widetilde g_{ij})\in C^1(\mathcal U,\Omega^1_{\UX/S})$, and therefore $\nabla(D)x=0$.

\noindent\textbf{2.10.} Functoriality of the crystalline Chern class implies that, for every $W$-valued point $e$ of $S$, there is a commutative square
\begin{equation}\tag{2.10.1}
\begin{tikzcd}[column sep=large,row sep=large]
\operatorname{Pic}(\UX_0)\arrow[r,"c_1"]\arrow[d,"\UL\mapsto \UL|_{X_0}"']&
H^2_{\mathrm{DR}}(\UX/S)=H^2_{\mathrm{cris}}(\UX_0/S)\arrow[d,"e^*"]\\
\operatorname{Pic}(X_0)\arrow[r,"c_1"']&H^2_{\mathrm{DR}}(X/W)=H^2_{\mathrm{cris}}(X_0/W),
\end{tikzcd}
\end{equation}
where $X/W$ is obtained from $\UX/S$ by base change along $e$. Since $X_0$ is a K3 surface, $\operatorname{Pic}(X_0)$ coincides with the Néron--Severi group $\operatorname{NS}(X_0)$ and is torsion free; hence the lower horizontal arrow in (2.10.1) is injective (3.4).
\newpage
\setcounter{footnote}{0}
\fontsize{10.2pt}{12.04pt}\selectfont

\noindent\textbf{2.11. Proof of 1.6.} Suppose that $L_0$ extends to an invertible sheaf $\UL_0$ on $\UX_0$, and let $x=c_1(\UL_0)\in H^2_{\mathrm{DR}}(\UX/S)$. Commutativity of (2.10.1) gives $e^*x=c_1(L_0)$. Since $L_0$ is nontrivial by hypothesis, $c_1(L_0)\ne0$ by the preceding remark, and hence $x\ne0$. By 2.9, $Fx=px$ and $\nabla x=0$. Let $p^n$ be the highest power of $p$ dividing $x$, and put $y=p^{-n}x$; the image $y_0$ of $y$ in
$H^2_{\mathrm{DR}}(\UX_0/S_0)=H^2_{\mathrm{DR}}(\UX/S)\bmod p$ is nonzero. We still have $Fy=py$ and $\nabla y=0$. By (2.6.1), the first relation implies
$y_0\in F^1_{\mathrm{Hdg}}H^2_{\mathrm{DR}}(\UX_0/S_0)$. Since $y_0\ne0$ and $\nabla y_0=0$, one cannot have
$y_0\in F^2_{\mathrm{Hdg}}H^2_{\mathrm{DR}}(\UX_0/S_0)$, for this would contradict the injectivity of
\[
\operatorname{gr}^2\nabla:F^2_{\mathrm{Hdg}}H^2_{\mathrm{DR}}(\UX_0/S_0)
\longrightarrow\Omega^1_{S_0/k}\otimes\operatorname{gr}^1_FH^2_{\mathrm{DR}}(\UX_0/S_0)
\]
from 2.4.4. Thus the image of $y_0$ in $\operatorname{gr}^1_FH^2_{\mathrm{DR}}(\UX_0/S_0)$ is nonzero; but $\nabla y_0=0$, contradicting the fact that
\[
\operatorname{gr}^1\nabla:\operatorname{gr}^1_FH^2_{\mathrm{DR}}(\UX_0/S_0)
\longrightarrow\Omega^1_{S_0/k}\otimes\operatorname{gr}^0_FH^2_{\mathrm{DR}}(\UX_0/S_0)
\]
is an isomorphism. This contradiction proves 1.6.

\section*{3. APPENDIX; CRYSTALLINE CHERN CLASSES AND INTERSECTIONS\footnote{By P. Deligne and L. Illusie.}}
The purpose of this appendix is to establish that, on a proper smooth surface over $k$, the intersection number of two divisors is the trace of the cup product of their crystalline Chern classes (this compatibility apparently does not occur in the literature).

\noindent\textbf{3.1.} Let $X$ be a proper smooth scheme over $k$. Recall [2, VI~3.3.6] that to every closed subscheme $Y$ of $X$, smooth over $k$ and of codimension $d$ in $X$, one can associate a cohomology class
\begin{equation}\tag{3.1.1}
c\ell(Y)\in H^{2d}(X/W).
\end{equation}
If $Y'$ and $Y''$ are smooth closed subschemes of codimensions $d'$ and $d''$, meeting transversely, then [2, VI~4.3.15]
\newpage
\setcounter{footnote}{0}
\fontsize{10.2pt}{12.04pt}\selectfont

\begin{equation}\tag{3.1.2}
c\ell(Y')\,c\ell(Y'')=c\ell(Y'\cap Y'')
\end{equation}
in $H^{2(d'+d'')}(X/W)$. If $X'$ is proper and smooth over $k$, $f:X'\to X$ is a $k$-morphism, and $Y$ is a smooth closed subscheme of codimension $d$ such that $Y\times_X X'$ is smooth of codimension $d$ in $X'$, then [2, VI~4.3.13]
\begin{equation}\tag{3.1.3}
c\ell(Y\times_X X')=f^*c\ell(Y).
\end{equation}
Finally, if $a$ is a rational point of $X/k$, then [2, VII~3.1.7]
\begin{equation}\tag{3.1.4}
\operatorname{Tr}_{X/W}(c\ell(a))=1.
\end{equation}
In particular, if $X$ is connected, $c\ell(a)$ is independent of $a$, since
$\operatorname{Tr}_{X/W}:H^{2n}(X/W)\to W$ is an isomorphism, where $n=\dim X$ [2, VII~2.1.1].

\noindent\textbf{3.2.} To every vector bundle $L$ on $X$ one can associate Chern classes [3, 2.4]
\begin{equation}\tag{3.2.1}
c_i(L)\in H^{2i}(X/W),
\end{equation}
which vanish for $i>\operatorname{rk}(L)$. These classes depend functorially on $L$: if $f:X'\to X$ is a $k$-morphism, then (loc.~cit.)
\begin{equation}\tag{3.2.2}
c_i(f^*L)=f^*c_i(L).
\end{equation}
If $D$ is a smooth effective divisor on $X$, there are therefore both the cohomology class $c\ell(D)\in H^2(X/W)$ and the Chern class $c_1(\mathcal O_X(D))\in H^2(X/W)$.

\noindent\textbf{PROPOSITION 3.2.3.} Let $D$ be a smooth effective divisor on $X$. If $\mathcal O_X(D)$ is very ample, then
\begin{equation}\tag{3.2.3.1}
c\ell(D)=c_1(\mathcal O_X(D)).
\end{equation}

\noindent\textbf{REMARK 3.2.4.} Berthelot (unpublished) verified that (3.2.3.1) holds without the hypothesis that $\mathcal O_X(D)$ be very ample. The proof is rather complicated. For our purposes, 3.2.3 is sufficient.
\newpage
\setcounter{footnote}{0}
\fontsize{10.2pt}{12.04pt}\selectfont

Let us prove 3.2.3. By 3.1.3 and 3.2.2, we reduce at once to the case in which $X$ is projective space $\mathbf P^N_k$ and $D$ is a hyperplane. Put $\mathbf P^N_W=X'$, and let $D'$ be a hyperplane of $X'$ lifting $D$. Then $H^2(X/W)=H^2_{\mathrm{DR}}(X'/W)$. By [2, VI~3.3.5], $c\ell(D)$ is the de Rham cohomology class of $D'$, which is computed from the extension of $\Omega^\bullet_{D'/W}[-1]$ by $\Omega^\bullet_{X'/W}$ supplied by the de Rham complex of $X'/W$ with logarithmic poles along $D'$. If $(t_0,\ldots,t_N)$ are homogeneous coordinates on $X'$, it follows readily that $c\ell(D)$ is represented, on the standard covering, by the $1$-cocycle
\[
\left(\frac{dt_j}{t_j}-\frac{dt_i}{t_i}\right)_{i<j}.
\]
On the other hand, by [3, 2.3], the Chern class $c_1(\mathcal O_{X'}(1))\in H^2_{\mathrm{DR}}(X'/W)$ is obtained from the cocycle $t_j/t_i$ by the homomorphism
$d\log:\mathcal O_{X'}^*\to\Omega^\bullet_{X'/W}[1]$. Thus $c_1(\mathcal O_X(D))=c\ell(D)$.

\noindent\textbf{THEOREM 3.3.} Let $X$ be a proper smooth surface over $k$. If $D_1,D_2$ are divisors on $X$, then
\[
(D_1\cdot D_2)=\operatorname{Tr}_{X/W}\bigl(c_1(\mathcal O_X(D_1))\,c_1(\mathcal O_X(D_2))\bigr),
\]
where $(D_1\cdot D_2)$ denotes their intersection number.

If $D_1,D_2$ are effective, smooth, very ample, and transverse, the conclusion follows from (3.1.2), (3.1.4), and 3.2.3. To reduce to this case, choose a very ample divisor $H$ on $X$. For $n$ sufficiently large, $D_1+nH$ and $D_2+nH$ are very ample. By Bertini, there are effective smooth very ample divisors $D'_1,D''_1,D'_2,D''_2$ such that
\[
\mathcal O_X(D_1)=\mathcal O_X(D'_1-D''_1),\qquad
\mathcal O_X(D_2)=\mathcal O_X(D'_2-D''_2),
\]
and one may further suppose that the four pairs
\[
(D'_1,D'_2),\quad(D'_1,D''_2),\quad(D''_1,D'_2),\quad(D''_1,D''_2)
\]
are transverse. Additivity of $c_1$ and bilinearity of the intersection pairing reduce the theorem to the special case already considered, completing the proof.
\newpage
\setcounter{footnote}{0}
\fontsize{10.2pt}{12.04pt}\selectfont

\noindent\textbf{3.4.} Again let $X$ be a proper smooth scheme over $k$. It is not known whether, for two smooth closed subschemes of $X$ of the same codimension, algebraic equivalence implies equality of their crystalline cohomology classes (this is known only for zero-dimensional subschemes; cf. 3.1.4!). On the other hand, if
$\operatorname{NS}(X)=\operatorname{Pic}(X)/\operatorname{Pic}^0(X)$ is the Néron--Severi group, the homomorphism
\[
c_1:\operatorname{Pic}(X)\longrightarrow H^2(X/W)
\]
vanishes on $\operatorname{Pic}^0(X)$, because $\operatorname{Pic}^0(X)$ is $p$-divisible and $H^2(X/W)$ is of finite type over $W$. It therefore induces
\begin{equation}\tag{3.4.1}
c_1:\operatorname{NS}(X)\longrightarrow H^2(X/W).
\end{equation}
Suppose now that $X$ is a surface. By the Hodge index theorem, the intersection form on $\operatorname{NS}(X)\otimes\mathbf Q$ is nondegenerate. In view of 3.3, it follows that
\begin{equation}\tag{3.4.2}
c_1\otimes\mathbf Q:\operatorname{NS}(X)\otimes\mathbf Q
\longrightarrow H^2(X/W)\otimes\mathbf Q
\end{equation}
is injective.

\noindent\textbf{REMARK 3.5.} One can show [6, II~5.10] that (3.4.1) defines an injection
$\operatorname{NS}(X)\otimes\mathbf Z_p\hookrightarrow H^2(X/W)$. It is probably not true in general that (3.4.1) induces an injection
$\operatorname{NS}(X)/p\operatorname{NS}(X)\to H^2(X/W)/pH^2(X/W)$, i.e. that the de Rham Chern-class map sends $\operatorname{NS}(X)/p\operatorname{NS}(X)$ injectively into $H^2_{\mathrm{DR}}(X/k)$. This does hold if $H^*(X/W)$ is torsion free and the Hodge-to-de Rham spectral sequence of $X/k$ degenerates at $E_1$; see [13, 1.4] and [6, II~5.18]. These conditions are satisfied, for example, by a K3 surface.

\vfill
\begin{minipage}[t]{0.46\textwidth}
\textbf{P. DELIGNE}\\
Institut des Hautes Etudes\\
Scientifiques\\
35, Route de Chartres\\
91440 BURES/YVETTE (France)
\end{minipage}\hfill
\begin{minipage}[t]{0.46\textwidth}
\textbf{L. ILLUSIE}\\
Université de Paris-Sud\\
Centre d'Orsay\\
Mathématique, bât.~425\\
91405 ORSAY (France)
\end{minipage}
\newpage
\setcounter{footnote}{0}
\fontsize{10.2pt}{12.04pt}\selectfont

\renewcommand{\refname}{\centering BIBLIOGRAPHY}
\begin{thebibliography}{16}\setlength{\itemsep}{0.35em}
\par\bibitem{1} M. ARTIN.- Supersingular K3 surfaces. Ann. Sc. ENS, 4ème série, t.~7 (1974), p.~543-568.
\par\bibitem{2} P. BERTHELOT.- Cohomologie cristalline des schémas de caractéristique $p>0$ . Lecture Notes in Math. 407, Springer-Verlag (1974).
\par\bibitem{3} P. BERTHELOT et L. ILLUSIE.- Classes de Chern en cohomologie cristalline. C. R. Acad. Sc. Paris, t.~270, p.~1695-1697 et 1750-1752, 22 et 29 juin 1970.
\par\bibitem{4} P. BERTHELOT et A. OGUS.- Notes on crystalline cohomology. Mathematical Notes 21, Princeton U. Press (1978).
\par\bibitem{5} P. DELIGNE.- La conjecture de Weil pour les surfaces K3. Inv. math. 15, p.~206-226 (1972).
\par\bibitem{6} L. ILLUSIE.- Complexe de de Rham-Witt et cohomologie cristalline, Ann. Sc. ENS, 4e série, t.~12 (1979), p.~501-661.
\par\bibitem{7} N. KATZ.- Travaux de Dwork, Séminaire Bourbaki, exp.~409, Lecture Notes in Math. 383, Springer-Verlag (1973).
\par\bibitem{8} N. KATZ.- Algebraic solutions of differential equations, $p$-curvature and the Hodge filtration. Inv. math. 18, p.~1-118 (1972).
\par\bibitem{9} N. KATZ and T. ODA.- On the differentiation of De Rham cohomology classes with respect to parameters. J. of Math. of Kyoto Univ., vol.~8, n$^\circ$~2, p.~199-213 (1968).
\par\bibitem{10} B. MAZUR.- Frobenius and the Hodge filtration, estimates. Ann. of Math. 98, p.~58-95 (1973).
\par\bibitem{11} D. MUMFORD.- Abelian Varieties. Tata Inst., Oxford Univ. Press (1970).
\par\bibitem{12} N. NYGAARD.- A $p$-adic proof of the Rudakov-Shafarevitch theorem. Ann. of Math. 110, p.~515-528 (1979).
\par\bibitem{13} A. OGUS.- Supersingular K3 crystals. Journées de Géométrie Algébrique de Rennes, juillet 78, S.M.F. Astérisque 64, p.~3-86 (1979).
\par\bibitem{14} A. RUDAKOV and I. SHAFAREVITCH.- Inseparable morphisms of algebraic surfaces. Akad. Sc. SSSR, t.~40, n$^\circ$~6, p.~1264-1307 (1976).
\par\bibitem{15} M. SCHLESSINGER.- Functors of Artin rings. Trans. Amer. Soc. 130, p.~205-222 (1968).
\par\bibitem{16} A. WEIL.- Variétés abéliennes et courbes algébriques. Hermann (1948).
\end{thebibliography}
\vfill
\end{document}
